\textbf{Phần 1. Trắc nghiệm lựa chọn}
\setcounterpageref{ex}{0}

\begin{ex} %[5H12N71]
Công thức tính diện tích hình tam giác có độ dài đáy $a$ và chiều cao $h$ (cùng đơn vị đo) là
\choice
{$S = a \times h$}
{$S = (a + h) \times 2$}
{$S = a + h : 2$}
{\True $S = \dfrac{a \times h}{2}$}
\loigiai{
$S = \dfrac{a \times h}{2}$.\\
\textbf{Đáp án của câu hỏi là D}
}
\end{ex}

\begin{ex} %[5H12N71]
Công thức tính diện tích hình thang có hai đáy $a, b$ và chiều cao $h$ là
\choice
{$S = (a + b) \times h$}
{$S = a \times b \times h$}
{\True $S = \dfrac{(a + b) \times h}{2}$}
{$S = (a + b) : 2 + h$}
\loigiai{
$S = \dfrac{(a + b) \times h}{2}$.\\
\textbf{Đáp án của câu hỏi là C}
}
\end{ex}

\begin{ex} %[5H12H71]
Một hình tròn có bán kính $5\text{ cm}$. Chu vi của hình tròn đó là
\choice
{\True $31{,}4\text{ cm}$}
{$15{,}7\text{ cm}$}
{$78{,}5\text{ cm}$}
{$62{,}8\text{ cm}$}
\loigiai{
$C = 2 \times 3{,}14 \times 5 = 31{,}4\text{ (cm)}$.\\
\textbf{Đáp án của câu hỏi là A}
}
\end{ex}

\begin{ex} %[5H12H71]
Một hình tròn có đường kính $6\text{ dm}$. Diện tích của hình tròn đó là
\choice
{$113{,}04\text{ dm}^2$}
{\True $28{,}26\text{ dm}^2$}
{$18{,}84\text{ dm}^2$}
{$56{,}52\text{ dm}^2$}
\loigiai{
Bán kính: $r = 3\text{ dm}$.\\
$S = 3{,}14 \times 3 \times 3 = 28{,}26\text{ (dm}^2\text{)}$.\\
\textbf{Đáp án của câu hỏi là B}
}
\end{ex}

\begin{ex} %[5H12H71]
Một hình thang có đáy lớn $15\text{ cm}$, đáy bé $9\text{ cm}$, chiều cao $8\text{ cm}$. Diện tích hình thang là
\choice
{\True $96\text{ cm}^2$}
{$192\text{ cm}^2$}
{$120\text{ cm}^2$}
{$84\text{ cm}^2$}
\loigiai{
$S = \dfrac{(15 + 9) \times 8}{2} = 96\text{ (cm}^2\text{)}$.\\
\textbf{Đáp án của câu hỏi là A}
}
\end{ex}

\begin{ex} %[5H12H71]
Một hình lập phương có cạnh $6\text{ cm}$. Diện tích toàn phần của nó là
\choice
{$144\text{ cm}^2$}
{\True $216\text{ cm}^2$}
{$216\text{ cm}^3$}
{$96\text{ cm}^2$}
\loigiai{
$S_{tp} = 6 \times 6 \times 6 = 216\text{ (cm}^2\text{)}$.\\
\textbf{Đáp án của câu hỏi là B}
}
\end{ex}

\begin{ex} %[5H12V71]
Một hình tam giác có diện tích $36\text{ cm}^2$, độ dài đáy là $12\text{ cm}$. Chiều cao tương ứng của tam giác là
\choice
{$3\text{ cm}$}
{$9\text{ cm}$}
{$12\text{ cm}$}
{\True $6\text{ cm}$}
\loigiai{
Chiều cao: $36 \times 2 : 12 = 6\text{ (cm)}$.\\
\textbf{Đáp án của câu hỏi là D}
}
\end{ex}

\begin{ex} %[5H12V71]
Một hình hộp chữ nhật có kích thước dài $8\text{ cm}$, rộng $6\text{ cm}$, cao $5\text{ cm}$. Thể tích của hình hộp chữ nhật đó là
\choice
{\True $240\text{ cm}^3$}
{$120\text{ cm}^3$}
{$140\text{ cm}^3$}
{$480\text{ cm}^3$}
\loigiai{
$V = 8 \times 6 \times 5 = 240\text{ (cm}^3\text{)}$.\\
\textbf{Đáp án của câu hỏi là A}
}
\end{ex}

\begin{ex} %[5H12V71]
Người ta quét sơn toàn bộ $6$ mặt ngoài của một khối gỗ hình lập phương có cạnh $1{,}2\text{ m}$. Diện tích quét sơn là
\choice
{$5{,}76\text{ m}^2$}
{$7{,}2\text{ m}^2$}
{\True $8{,}64\text{ m}^2$}
{$1{,}728\text{ m}^2$}
\loigiai{
$S_{tp} = 1{,}2 \times 1{,}2 \times 6 = 8{,}64\text{ (m}^2\text{)}$.\\
\textbf{Đáp án của câu hỏi là C}
}
\end{ex}

\begin{ex} %[5H12V71]
Một mảnh đất hình thang có diện tích $180\text{ m}^2$, chiều cao $12\text{ m}$. Trung bình cộng hai đáy của hình thang là
\choice
{$30\text{ m}$}
{\True $15\text{ m}$}
{$20\text{ m}$}
{$25\text{ m}$}
\loigiai{
Trung bình cộng 2 đáy: $180 : 12 = 15\text{ (m)}$.\\
\textbf{Đáp án của câu hỏi là B}
}
\end{ex}

\begin{ex} %[5H12V71]
Một bồn hoa gồm một hình vuông cạnh $4\text{ m}$ và $4$ nửa hình tròn có đường kính là cạnh của hình vuông gắn vào 4 cạnh. Diện tích toàn bộ bồn hoa là
\choice
{$28{,}56\text{ m}^2$}
{$32{,}56\text{ m}^2$}
{\True $41{,}12\text{ m}^2$}
{$50\text{ m}^2$}
\loigiai{
Diện tích hình vuông: $4 \times 4 = 16\text{ (m}^2\text{)}$.\\
4 nửa hình tròn = 2 hình tròn bán kính $2\text{ m}$.\\
Diện tích 2 hình tròn: $2 \times (3{,}14 \times 2 \times 2) = 25{,}12\text{ (m}^2\text{)}$.\\
Tổng diện tích: $16 + 25{,}12 = 41{,}12\text{ (m}^2\text{)}$.\\
\textbf{Đáp án của câu hỏi là C}
}
\end{ex}

\begin{ex} %[5H12C71]
Cho hình thang vuông $ABCD$ vuông tại $A$ và $D$, có đáy bé $AB = 15\text{ cm}$, đáy lớn $CD = 25\text{ cm}$, chiều cao $AD = 16\text{ cm}$. Điểm $M$ là trung điểm của cạnh $AD$. Diện tích tam giác $MBC$ là
\choice
{$140\text{ cm}^2$}
{$180\text{ cm}^2$}
{$200\text{ cm}^2$}
{\True $160\text{ cm}^2$}
\loigiai{
$S_{ABCD} = \dfrac{(15 + 25) \times 16}{2} = 320\text{ (cm}^2\text{)}$.\\
$M$ là trung điểm $AD \Rightarrow AM = MD = 8\text{ cm}$.\\
$S_{ABM} = \dfrac{15 \times 8}{2} = 60\text{ (cm}^2\text{)}$.\\
$S_{MCD} = \dfrac{25 \times 8}{2} = 100\text{ (cm}^2\text{)}$.\\
$S_{MBC} = 320 - (60 + 100) = 160\text{ (cm}^2\text{)}$.\\
\textbf{Đáp án của câu hỏi là D}
}
\end{ex}

\textbf{Phần 2. Trắc nghiệm đúng sai}
\setcounterpageref{ex}{0}

\begin{ex} %[5H12H71]
Cho tam giác $ABC$ có đáy $BC = 16\text{ cm}$, chiều cao $AH = 10\text{ cm}$. Xét tính đúng sai:
\choiceTF
{\True Diện tích tam giác $ABC$ là $80\text{ cm}^2$.}
{Nếu tăng đáy $BC$ thêm $4\text{ cm}$ thì diện tích tam giác tăng thêm $40\text{ cm}^2$.}
{\True Nếu tăng đáy $BC$ thêm $4\text{ cm}$ thì diện tích tam giác tăng thêm $20\text{ cm}^2$.}
{\True Khi chiều cao tăng gấp đôi thì diện tích tam giác tăng gấp đôi.}
\loigiai{
\begin{itemize}
    \item $S = 16 \times 10 : 2 = 80\text{ cm}^2$. Đúng.
    \item Diện tích tăng thêm: $4 \times 10 : 2 = 20\text{ cm}^2 \neq 40\text{ cm}^2$. Ý 2 Sai, Ý 3 Đúng.
    \item Diện tích tỉ lệ thuận với chiều cao nên tăng gấp đôi. Đúng.
\end{itemize}
\begin{itemize}
    \item a) \textbf{Đúng}.
    \item b) \textbf{Sai}.
    \item c) \textbf{Đúng}.
    \item d) \textbf{Đúng}.
\end{itemize}
}
\end{ex}

\begin{ex} %[5H12V71]
Một bể nước hình hộp chữ nhật có kích thước lòng trong: dài $2\text{ m}$, rộng $1{,}5\text{ m}$, cao $1{,}2\text{ m}$. Xét tính đúng sai:
\choiceTF
{\True Diện tích đáy bể là $3\text{ m}^2$.}
{\True Diện tích xung quanh bên trong bể là $8{,}4\text{ m}^2$.}
{Dung tích tối đa của bể là $3\,000\text{ l}$.}
{\True Dung tích tối đa của bể là $3\,600\text{ l}$.}
\loigiai{
\begin{itemize}
    \item $S_{đáy} = 2 \times 1{,}5 = 3\text{ m}^2$. Đúng.
    \item $S_{xq} = (2 + 1{,}5) \times 2 \times 1{,}2 = 8{,}4\text{ m}^2$. Đúng.
    \item Dung tích: $3 \times 1{,}2 = 3{,}6\text{ m}^3 = 3600\text{ l} \neq 3000\text{ l}$. Ý 3 Sai, Ý 4 Đúng.
\end{itemize}
\begin{itemize}
    \item a) \textbf{Đúng}.
    \item b) \textbf{Đúng}.
    \item c) \textbf{Sai}.
    \item d) \textbf{Đúng}.
\end{itemize}
}
\end{ex}

\textbf{Phần 3. Trả lời ngắn}
\setcounterpageref{ex}{0}

\begin{ex} %[5H12H71]
Một hình tam giác có đáy $14\text{ cm}$, chiều cao $8\text{ cm}$. Tính diện tích theo xăng-ti-mét vuông.
\shortans{56}
\loigiai{
\textbf{Đáp số:} $56$\\
$14 \times 8 : 2 = 56\text{ cm}^2$.
}
\end{ex}

\begin{ex} %[5H12H71]
Một hình tròn có bán kính $4\text{ cm}$. Tính diện tích theo xăng-ti-mét vuông.
\shortans{50,24}
\loigiai{
\textbf{Đáp số:} $50{,}24$\\
$3{,}14 \times 4 \times 4 = 50{,}24\text{ cm}^2$.
}
\end{ex}

\begin{ex} %[5H12V71]
Một hình thang có đáy lớn $20\text{ m}$, đáy bé $14\text{ m}$, chiều cao $10\text{ m}$. Tính diện tích hình thang theo mét vuông.
\shortans{170}
\loigiai{
\textbf{Đáp số:} $170$\\
$(20 + 14) \times 10 : 2 = 170\text{ m}^2$.
}
\end{ex}

\begin{ex} %[5H12V71]
Một hình lập phương có diện tích toàn phần là $96\text{ cm}^2$. Tính thể tích hình lập phương theo xăng-ti-mét khối.
\shortans{64}
\loigiai{
\textbf{Đáp số:} $64$\\
Diện tích 1 mặt: $96 : 6 = 16\text{ cm}^2 \Rightarrow$ Cạnh $4\text{ cm}$.\\
$V = 4 \times 4 \times 4 = 64\text{ cm}^3$.
}
\end{ex}

\begin{ex} %[5H12V71]
Một bể cá hình hộp chữ nhật dài $80\text{ cm}$, rộng $50\text{ cm}$, cao $40\text{ cm}$. Tính dung tích tối đa của bể theo lít.
\shortans{160}
\loigiai{
\textbf{Đáp số:} $160$\\
$8 \times 5 \times 4 = 160\text{ dm}^3 = 160\text{ l}$.
}
\end{ex}

\begin{ex} %[5H12C71]
Một cái sân hình chữ nhật dài $18\text{ m}$, rộng $12\text{ m}$. Ở giữa sân người ta xây một bồn hoa hình tròn bán kính $2\text{ m}$. Tính diện tích phần sân còn lại theo mét vuông.
\shortans{203,44}
\loigiai{
\textbf{Đáp số:} $203{,}44$\\
Diện tích sân: $18 \times 12 = 216\text{ m}^2$.\\
Diện tích bồn hoa: $3{,}14 \times 2 \times 2 = 12{,}56\text{ m}^2$.\\
Diện tích còn lại: $216 - 12{,}56 = 203{,}44\text{ m}^2$.
}
\end{ex}
